\subsection{MENTOR行业排序}
模型输入文本覆盖2023--2024年，包括160,190篇中文KOL帖子与72,108篇英文财经新闻，A股和美股行情分别来自Wind与Bloomberg。排序评价涉及9个与11个行业。\src{FUND-113}\src{FUND-114}

\begin{table}[H]\small\centering
\caption{MENTOR论文中的行业排序结果}\label{tab:mentor}
\begin{tabularx}{\linewidth}{@{}P{41mm}Y Y Y@{}}\toprule
方法 & 美股行业排序相关性 & A股行业排序相关性 & A股前3行业命中率\\\midrule
动量基线 & 0.159 & 0.100 & 0.444\\
SEP + DeepSeek & 0.078 & 0.118 & 0.472\\
MENTOR + DeepSeek & 0.220 & 0.141 & 0.488\\
MENTOR + O1 & 0.221 & 0.173 & 0.439\\\bottomrule
\end{tabularx}
\source{MENTOR正文表3（PDF第11页）。排序相关性采用Spearman系数；前3命中率衡量预测与实际领先行业的重合程度。数字为原作者报告值。}
\end{table}
表中MENTOR改善了部分排序相关性，但A股O1版本的前3命中率低于DeepSeek版本与动量。完整版权数据未公开，正文和补充材料部分数据描述对应不清，预训练污染与样本时间仍需复核。

\subsection{Man AlphaTrend验证实验}
作者故意从既有突破信号中移除提高反应速度的机制，再输入预期有效、预期较差和开放性想法，检查研究流程能否形成相应结论。\src{FUND-323}

\begin{table}[H]\small\centering
\caption{AlphaTrend对三类研究想法的检验}\label{tab:man}
\begin{tabularx}{\linewidth}{@{}P{38mm}Y Y@{}}\toprule
研究命题 & 原作者观察到的结果 & 结果的含义\\\midrule
恢复被移除的反应机制 & 方案的Sharpe集中在约1.05，基准约为0.95。 & 系统识别出预先设定的改进机会。\\
改用局部峰谷等信号 & Sharpe中心约为0.80，几乎所有方案落后于基准。 & 系统能够给出负面结果，帮助停止无效尝试。\\
在更宽的范围探索趋势方案 & Sharpe约为0.70--1.00，中心约0.85，不同时期表现不同。 & 结果存在条件差异，需要进一步限定假设与适用范围。\\\bottomrule
\end{tabularx}
\source{Man AHL，2026-02-11，图2--4及相邻正文。结果是原作者的历史模拟展示，不代表真实产品业绩。}
\end{table}
恢复已知机制主要验证流程能否工作，不代表发现新Alpha。原文未完整披露资产池、参数、成本和统计检验；图注截至2015年，正文部分表述涉及2012--2016年，不能补写统一样本区间。同一开放想法下，作者报告Claude 4.0 Sonnet方案的大部分两两相关性高于0.85，GPT-5约为0.75--1.0，仅说明该设置中的候选差异。

\subsection{QuantaAlpha总体实验及案例口径}
表中使用同一GPT-5.2模型比较三个LLM框架，各方法产生约150个通过验证的因子，再交给同一LightGBM模型评价。\src{EXT-QA}\src{EXT-QC}

\begin{table}[H]\small\centering
\caption{QuantaAlpha论文中的同模型对照}\label{tab:qa}
\begin{tabularx}{\linewidth}{@{}P{43mm}Y Y Y@{}}\toprule
框架（均用GPT-5.2） & IC & 年化超额收益 & 超额收益最大回撤\\\midrule
RD-Agent & 0.0286 & 3.58\% & 16.76\%\\
AlphaAgent & 0.0347 & 1.11\% & 13.89\%\\
QuantaAlpha & 0.0472 & 4.68\% & 11.80\%\\\bottomrule
\end{tabularx}
\source{QuantaAlpha v3表1及附录A--B。IC衡量预测值与后续收益的相关性，表中数字均为原作者报告值。}
\end{table}
总体市场为沪深300，训练期2016--2020年，验证期2021年，测试期2022年初至2025年12月26日。组合每日选择50只等权股票，每次替换5只，采用次日开盘价，买入与卖出费率分别为0.05\%和0.15\%。收益与回撤基于扣费后的超额收益；预测标签为下一交易日到再下一交易日的收盘收益。其他传统或深度方法在部分组合指标上更好，不能据该节选认定LLM全面占优。

案例表达式的价格与成交量变动均以当日值作分母，日内收益也以当日收盘价作分母。复现时应按原表达式保留这一口径。

使用DeepSeek-V3.2的消融实验中，移除修改环节后，IC由0.0461降至0.0382，年化超额收益由4.53\%降至3.27\%。附录C案例则报告子因子Rank IC为0.0311，父记录为0.0216和0.0246，信息比率由0.973降至0.963，超额收益最大回撤由约7.30\%扩大至11.37\%。该案例与总体实验的轮次、换手及成本口径不完全一致，正文仅用其说明模型解释反馈和提出修改的过程，不合并为总体收益证据。

作者报告一次主要运行约20小时、约180万tokens，采用托管模型API与CPU回测。公开代码为选取文件存档，配置、数据许可和完整运行环境仍需核对，尚未执行复现。

\subsection{AlphaQT-Bench代码生成评价}
该基准包含270项任务并测试12个模型，LLM返回Python实现，测试程序分别检查运行、时序因果、功能和计算结构。四项同时通过的平均验证正确率（Verified Accuracy），在论文模型样本中最高为93.7\%（Gemini-2.5-Pro），开放权重模型最高为81.9\%（DeepSeek-V3）。该结果评价任务实现，不等于策略成功率或本地推理性能。\src{EXT-AB}
